% Options for packages loaded elsewhere
\PassOptionsToPackage{unicode}{hyperref}
\PassOptionsToPackage{hyphens}{url}
\PassOptionsToPackage{dvipsnames,svgnames,x11names}{xcolor}
\documentclass[
  english,
  10pt,
  a4paper,
]{article}
\usepackage{xcolor}
\usepackage[top=22mm,bottom=20mm,left=18mm,right=18mm]{geometry}
\usepackage{amsmath,amssymb}
\setcounter{secnumdepth}{-\maxdimen} % remove section numbering
\usepackage{iftex}
\ifPDFTeX
  \usepackage[T1]{fontenc}
  \usepackage[utf8]{inputenc}
  \usepackage{textcomp} % provide euro and other symbols
\else % if luatex or xetex
  \usepackage{unicode-math} % this also loads fontspec
  \defaultfontfeatures{Scale=MatchLowercase}
  \defaultfontfeatures[\rmfamily]{Ligatures=TeX,Scale=1}
\fi
\usepackage{lmodern}
\ifPDFTeX\else
  % xetex/luatex font selection
\fi
% Use upquote if available, for straight quotes in verbatim environments
\IfFileExists{upquote.sty}{\usepackage{upquote}}{}
\IfFileExists{microtype.sty}{% use microtype if available
  \usepackage[]{microtype}
  \UseMicrotypeSet[protrusion]{basicmath} % disable protrusion for tt fonts
}{}
\makeatletter
\@ifundefined{KOMAClassName}{% if non-KOMA class
  \IfFileExists{parskip.sty}{%
    \usepackage{parskip}
  }{% else
    \setlength{\parindent}{0pt}
    \setlength{\parskip}{6pt plus 2pt minus 1pt}}
}{% if KOMA class
  \KOMAoptions{parskip=half}}
\makeatother
\usepackage{longtable,booktabs,array}
\usepackage{caption}
\captionsetup[table]{skip=6pt}
\newcounter{none} % for unnumbered tables
\usepackage{calc} % for calculating minipage widths
% Correct order of tables after \paragraph or \subparagraph
\usepackage{etoolbox}
\makeatletter
\patchcmd\longtable{\par}{\if@noskipsec\mbox{}\fi\par}{}{}
\makeatother
\ifLuaTeX
\usepackage[bidi=basic,shorthands=off]{babel}
\else
\usepackage[bidi=default,shorthands=off]{babel}
\fi
\ifLuaTeX
  \usepackage{selnolig} % disable illegal ligatures
\fi
\setlength{\emergencystretch}{3em} % prevent overfull lines
\providecommand{\tightlist}{%
  \setlength{\itemsep}{0pt}\setlength{\parskip}{0pt}}
% ---- Çılgın Kurye document theme: same look as the CSE 396 project proposal ----
% Compile with XeLaTeX (also works on Overleaf). Fonts: TeX Gyre Termes (Times), DejaVu Sans Mono.
\usepackage{fontspec}
\setmainfont{texgyretermes}[Extension=.otf, UprightFont=*-regular, BoldFont=*-bold, ItalicFont=*-italic, BoldItalicFont=*-bolditalic]
\setsansfont{texgyreheros}[Extension=.otf, UprightFont=*-regular, BoldFont=*-bold, ItalicFont=*-italic, BoldItalicFont=*-bolditalic]
\setmonofont{DejaVuSansMono}[Extension=.ttf, UprightFont=*, BoldFont=*-Bold, ItalicFont=*-Oblique, BoldItalicFont=*-BoldOblique, Scale=0.80]
\setmathfont{texgyretermes-math.otf}

\usepackage{xcolor}
\definecolor{ckblue}{HTML}{1F3A5F}
\definecolor{cklight}{HTML}{E8EEF5}
\definecolor{ckgrey}{HTML}{F2F2F2}
\definecolor{ckline}{HTML}{555555}
\definecolor{shadecolor}{HTML}{F4F5F7}

% Headings as in the proposal: serif, bold; subsections bold italic. Numbers are part of the heading text.
\setcounter{secnumdepth}{0}
\usepackage{titlesec}
\titleformat{\section}{\rmfamily\Large\bfseries}{}{0pt}{}
\titleformat{\subsection}{\rmfamily\large\bfseries\itshape}{}{0pt}{}
\titleformat{\subsubsection}{\rmfamily\normalsize\bfseries}{}{0pt}{}
\titlespacing*{\section}{0pt}{12pt}{4pt}
\titlespacing*{\subsection}{0pt}{8pt}{3pt}

% Header and footer (the first page has only the footer, like the proposal)
\usepackage{fancyhdr,lastpage}
\pagestyle{fancy}
\fancyhf{}
\fancyhead[L]{\footnotesize\itshape CSE 396 · Group 8 · Çılgın Kurye}
\fancyhead[R]{\footnotesize\itshape\ckmodule}
\fancyfoot[C]{\footnotesize Page \thepage\ of \pageref*{LastPage}}
\renewcommand{\headrulewidth}{0.4pt}
\setlength{\headheight}{14pt}
\fancypagestyle{ckfirst}{\fancyhf{}\renewcommand{\headrulewidth}{0pt}\fancyfoot[C]{\footnotesize Page \thepage\ of \pageref*{LastPage}}}

% Inline code: breaks only after . : ( , / _ and spaces, never inside a name
\usepackage[obeyspaces,spaces]{url}
\AtBeginDocument{%
  \renewcommand{\UrlBreaks}{\do\.\do\:\do\(\do\,\do\/\do\_\do\|}%
  \renewcommand{\UrlBigBreaks}{}%
  \renewcommand{\UrlNoBreaks}{}%
  \urlstyle{tt}}

% Tables: boxed grid as in the proposal, smaller text
\usepackage{etoolbox}
\AtBeginEnvironment{longtable}{\footnotesize\renewcommand{\arraystretch}{1.2}}
\setlength{\LTpre}{6pt}
\setlength{\LTpost}{8pt}
\setlength{\arrayrulewidth}{0.4pt}

% Code blocks: long lines wrap; diagrams stay in one piece
\usepackage{framed}
\usepackage{fvextra}
\RecustomVerbatimEnvironment{verbatim}{Verbatim}{fontsize=\footnotesize, breaklines}
\makeatletter
\@ifundefined{Shaded}{}{\renewenvironment{Shaded}{\begin{snugshade}}{\end{snugshade}}}
\@ifundefined{Highlighting}{}{\RecustomVerbatimEnvironment{Highlighting}{Verbatim}{commandchars=\\\{\}, fontsize=\footnotesize, breaklines, breakanywhere}}
\makeatother

% Figures drawn with TikZ
\usepackage{tikz}
\usetikzlibrary{arrows.meta,positioning,fit,calc,backgrounds,shapes.geometric}
\tikzset{
  ck/.style={font=\footnotesize, >={Stealth[length=5pt]}, line width=0.6pt},
  box/.style={draw=ckline, fill=white, rounded corners=2pt, align=center, inner sep=4pt, minimum height=7mm},
  mod/.style={box, fill=ckgrey},
  this/.style={box, fill=cklight, draw=ckblue, line width=0.9pt},
  group/.style={draw=ckline, dashed, rounded corners=3pt, inner sep=6pt},
  lbl/.style={font=\scriptsize, align=center, inner sep=1.5pt},
  msg/.style={->, draw=ckline},
}
\usepackage{caption}
\captionsetup[figure]{font=small, labelfont=bf, skip=4pt}
\newcommand{\ckfigure}[2]{\begin{figure}[htbp]\centering #1\caption{#2}\end{figure}}

\usepackage{enumitem}
\setlist{itemsep=1pt, topsep=3pt}
\setlength{\parskip}{4pt}
\setlength{\parindent}{0pt}
\raggedbottom
% ---- end of theme ----
\newcommand{\ckmodule}{M4 Traffic and AI}
\usepackage{bookmark}
\IfFileExists{xurl.sty}{\usepackage{xurl}}{} % add URL line breaks if available
\urlstyle{same}
% fallback for those not using the hyperref driver hyperxmp:
\makeatletter
\@ifundefined{xmpquote}{\newcommand{\xmpquote}[1]{#1}}{}
\makeatother
\hypersetup{
  pdftitle={M4 Traffic and AI},
  pdflang={en},
  colorlinks=true,
  linkcolor={black},
  filecolor={Maroon},
  citecolor={Blue},
  urlcolor={black},
  pdfcreator={LaTeX}}

\title{M4 Traffic and AI --- Module Documentation}
\author{}
\date{}


% ---- colours and drawings ----
\usepackage{colortbl}
\usepackage{float}
% ---- colours ----
\definecolor{srv}{HTML}{2E6FB7}\colorlet{srvL}{srv!15}
\definecolor{cli}{HTML}{2E9E5B}\colorlet{cliL}{cli!15}
\definecolor{adm}{HTML}{E08A1E}\colorlet{admL}{adm!18}
\definecolor{brd}{HTML}{7B4FB3}\colorlet{brdL}{brd!15}
\definecolor{tcpc}{HTML}{1F5FA8}
\definecolor{udpc}{HTML}{D9531E}
\definecolor{obs}{HTML}{D7263D}
\definecolor{tA}{HTML}{00A6D6}
\definecolor{tB}{HTML}{3BB273}
\definecolor{tC}{HTML}{F2C14E}
\definecolor{tD}{HTML}{9B5DE5}

% ---- styles ----
\tikzset{
  nb/.style={rounded corners=3pt, line width=0.9pt, align=center, inner sep=4pt, minimum height=9mm, font=\footnotesize},
  sv/.style={nb, draw=srv, fill=srvL},
  cl/.style={nb, draw=cli, fill=cliL},
  ad/.style={nb, draw=adm, fill=admL},
  bd/.style={nb, draw=brd, fill=brdL},
  hwb/.style={nb, draw=black!50, fill=black!5},
  tcpa/.style={draw=tcpc, line width=1.5pt},
  udpa/.style={draw=udpc, line width=1.5pt, dash pattern=on 5pt off 2.5pt},
  inna/.style={draw=black!55, line width=1pt},
  al/.style={font=\scriptsize, fill=white, inner sep=1.2pt, rounded corners=1pt},
  grp/.style={draw=black!35, dashed, rounded corners=5pt, inner sep=7pt},
  glbl/.style={font=\scriptsize\bfseries, text=black!65},
  hl/.style={line width=2.6pt},
  st/.style={nb, rounded corners=8pt, minimum width=24mm},
  co/.style={draw=black!60, line width=0.5pt},
  cot/.style={font=\scriptsize, align=left, text=black!80},
}

% ---- system map; argument = module to highlight (M1 ... M8, or none) ----
\newcommand{\systemmap}[1]{%
\tikzset{M1/.style={},M2/.style={},M3/.style={},M4/.style={},M5/.style={},M6/.style={},M7/.style={},M8/.style={},none/.style={}}%
\tikzset{#1/.style={hl}}%
\begin{tikzpicture}[ck]
\node[cl, M5, minimum width=27mm] (m5) at (0,3) {\textbf{M5}\\Player phones};
\node[ad, M6, minimum width=27mm] (m6) at (0,0.4) {\textbf{M6}\\Admin panel};
\node[sv, M2, minimum width=27mm, minimum height=38mm] (m2) at (4.6,1.7) {\textbf{M2}\\Network\\and sync};
\node[sv, M1, minimum width=27mm, minimum height=38mm] (m1) at (8.4,1.7) {\textbf{M1}\\Game server\\and rules};
\node[sv, M3, minimum width=25mm] (m3) at (12,2.8) {\textbf{M3}\\Physics};
\node[sv, M4, minimum width=25mm] (m4) at (12,0.6) {\textbf{M4}\\Traffic and AI};
\node[grp, fit=(m2)(m1)(m3)(m4)] (pc) {};
\node[glbl, anchor=south west] at (pc.north west) {MAIN COMPUTER};
\node[bd, M7, minimum width=27mm] (m7) at (4.6,-2.7) {\textbf{M7}\\ESP32 software};
\node[bd, M8, minimum width=27mm] (m8) at (8.4,-2.7) {\textbf{M8}\\LED engine};
\node[hwb, minimum width=25mm] (led) at (12,-2.7) {RGB LED\\strips};
\node[hwb, minimum width=27mm] (sm) at (4.6,-4.4) {Sensor and\\4 motors};
\node[grp, fit=(m7)(m8)(led)(sm)] (bb) {};
\node[glbl, anchor=north east] at (bb.south east) {LED CITY BOARD};
\draw[tcpa, <->] (m5.east) -- node[al]{WebSocket} (m5.east -| m2.west);
\draw[tcpa, <->] (m6.east) -- node[al]{WebSocket} (m6.east -| m2.west);
\draw[tcpa, <->] ([xshift=-6mm]m7.north) -- node[al]{TCP} ([xshift=-6mm]m2.south);
\draw[udpa, <->] ([xshift=6mm]m7.north) -- node[al]{UDP} ([xshift=6mm]m2.south);
\draw[inna, <->] (m2) -- (m1);
\draw[inna, <->] (m1.east |- m3) -- (m3);
\draw[inna, <->] (m1.east |- m4) -- (m4);
\draw[inna, ->] (m7) -- node[al]{state} (m8);
\draw[inna, ->] (m8) -- (led);
\draw[inna, <->] (m7) -- (sm);
\ckleg{-1.3,-6.0}{tcpa/TCP or WebSocket: every message arrives in order, udpa/UDP: fast and the newest message wins, inna/inside one program or chip}
\end{tikzpicture}}

% ---- LED board seen from above (6 x 6 roads, 1 unit = 10 cm) ----
\newcommand{\boardbase}{%
\fill[black!88, rounded corners=3pt] (-0.5,-0.5) rectangle (5.5,5.5);
\foreach \i in {0,...,4}{\foreach \j in {0,...,4}{%
  \fill[black!60] (\i+0.13,\j+0.13) rectangle (\i+0.87,\j+0.87);
  \fill[black!22] (\i+0.24,\j+0.24) rectangle (\i+0.52,\j+0.52);
  \fill[black!30] (\i+0.56,\j+0.30) rectangle (\i+0.78,\j+0.62);
  \fill[black!38] (\i+0.28,\j+0.58) rectangle (\i+0.66,\j+0.78);}}
\foreach \i in {0,...,5}{%
  \draw[black, line width=4.5pt] (\i,-0.05) -- (\i,5.05);
  \draw[black, line width=4.5pt] (-0.05,\i) -- (5.05,\i);}
\foreach \i in {0,...,5}{\foreach \k in {0,...,30}{%
  \fill[black!55] (\i,\k/6) circle (0.025);
  \fill[black!55] (\k/6,\i) circle (0.025);}}
\foreach \x/\y/\c in {0/4/tA, 4/4/tB, 0/0/tD, 4/0/tC}{%
  \fill[\c!30!black!60] (\x+0.13,\y+0.13) rectangle (\x+0.87,\y+0.87);
  \fill[\c] (\x+0.3,\y+0.3) rectangle (\x+0.7,\y+0.7);
  \fill[white] (\x+0.4,\y+0.5) circle (0.05); \fill[white] (\x+0.6,\y+0.5) circle (0.05);}}
\newcommand{\nmark}[3]{\node[circle, fill=yellow!85!orange, draw=black, line width=0.4pt, inner sep=0.6pt, minimum size=3.6mm, font=\tiny\bfseries] at (#1,#2) {#3};}
\newcommand{\bike}[4]{%
\fill[#3!35] ($(#1,#2)+(#4:0.40)$) circle (0.05);
\fill[#3!65] ($(#1,#2)+(#4:0.21)$) circle (0.075);
\fill[#3, draw=white, line width=0.4pt] (#1,#2) circle (0.12);}
\newcommand{\dpoint}[3]{\draw[#3!50, line width=0.8pt] (#1,#2) circle (0.26); \draw[#3, line width=1.4pt] (#1,#2) circle (0.16); \fill[#3] (#1,#2) circle (0.07);}
\newcommand{\heavy}[2]{\draw[obs!60!black, line width=4pt] (#1) -- (#2); \draw[obs, line width=4pt, dash pattern=on 3pt off 5pt] (#1) -- (#2);}
\newcommand{\roadwork}[2]{\draw[obs, line width=4pt] (#1) -- (#2);}
\newcommand{\accident}[2]{\draw[obs, line width=4pt, dash pattern=on 3pt off 3pt] (#1) -- (#2);}
\newcommand{\puddle}[2]{\draw[cyan!55, line width=4pt] (#1) -- (#2); \draw[cyan!15, line width=1.6pt, dash pattern=on 1pt off 2.5pt] (#1) -- (#2);}
\newcommand{\findme}[2]{\fill[white, opacity=0.25] (#1,#2) circle (0.36); \fill[white, opacity=0.5] (#1,#2) circle (0.24); \fill[white] (#1,#2) circle (0.13);}

% ---- phone game screen (7 x 3.4 units) ----
\newcommand{\phonemock}{%
\fill[black!85, rounded corners=9pt] (-0.3,-0.3) rectangle (7.3,3.7);
\fill[cyan!20] (0,1.7) rectangle (7,3.4);
\fill[cli!14] (0,0) rectangle (7,1.7);
\filldraw[fill=srv!30, draw=srv!70!black] (0.2,1.5) rectangle (1.0,3.0);
\filldraw[fill=adm!35, draw=adm!70!black] (1.1,1.6) rectangle (1.8,2.4);
\filldraw[fill=brd!30, draw=brd!70!black] (1.9,1.65) -- (2.6,1.65) -- (2.6,2.2) -- (2.25,2.5) -- (1.9,2.2) -- cycle;
\filldraw[fill=adm!35, draw=adm!70!black] (4.5,1.65) rectangle (5.1,2.3);
\filldraw[fill=srv!30, draw=srv!70!black] (5.2,1.55) rectangle (6.0,3.1);
\filldraw[fill=cli!35, draw=cli!70!black] (6.1,1.6) rectangle (6.85,2.6);
\fill[brown!60!black] (4.17,1.6) rectangle (4.23,1.8); \fill[cli!70!black] (4.2,1.95) circle (0.18);
\fill[black!55] (1.8,0) -- (5.2,0) -- (3.75,1.7) -- (3.25,1.7) -- cycle;
\draw[white, line width=0.8pt, dash pattern=on 4pt off 4pt] (3.5,0) -- (3.5,1.7);
\filldraw[fill=black!15, draw=black!70] (3.62,1.12) rectangle (3.88,1.32);
\fill[obs] (3.05,1.22) -- (3.15,1.48) -- (3.25,1.22) -- cycle;
\fill[black!85] (3.43,0.08) rectangle (3.57,0.3);
\fill[tA] (3.3,0.28) rectangle (3.7,0.6);
\fill[brown!65] (3.34,0.6) rectangle (3.66,0.82);
\fill[tA!60!black] (3.5,0.95) circle (0.13);
\node[fill=black!70, text=white, rounded corners=2pt, font=\scriptsize\bfseries, inner sep=2pt] at (3.5,3.15) {2:41};
\foreach \c/\s/\x in {tA/120/0.35,tB/95/0.95,tC/140/1.55,tD/80/2.15}
  \node[fill=\c, text=white, font=\tiny\bfseries, rounded corners=1pt, inner sep=1.5pt, minimum width=5mm] at (\x,3.15) {\s};
\fill[white, opacity=0.9] (6.45,2.85) circle (0.38); \draw[black!60] (6.45,2.85) circle (0.38);
\draw[-{Stealth[length=5pt]}, line width=1.8pt, obs] (6.45,2.85) ++(230:0.2) -- ++(50:0.5);
\fill[white, opacity=0.85] (0.55,0.5) circle (0.32); \fill[black!70] (0.43,0.5) -- (0.65,0.63) -- (0.65,0.37) -- cycle;
\fill[white, opacity=0.85] (1.4,0.5) circle (0.32); \fill[black!70] (1.52,0.5) -- (1.3,0.63) -- (1.3,0.37) -- cycle;
\node[fill=obs!80, text=white, rounded corners=3pt, font=\tiny\bfseries, minimum width=8mm, minimum height=6mm] at (5.6,0.5) {Brake};
\node[fill=cli, text=white, rounded corners=3pt, font=\tiny\bfseries, minimum width=8mm, minimum height=8mm] at (6.5,0.6) {Gas};
\node[fill=tC, text=black, rounded corners=3pt, font=\tiny\bfseries, minimum width=9mm, minimum height=5mm] at (6.15,1.45) {Deliver};
\node[fill=white, text=black, rounded corners=3pt, font=\tiny\bfseries, minimum width=9mm, minimum height=5mm] at (0.9,1.45) {Find me};}

% ---- side cut of the board (14 units wide) ----
\newcommand{\boardside}{%
\fill[brown!20] (-0.6,-0.35) rectangle (14.6,0);
\foreach \x/\h/\c in {1.0/0.7/black!18, 2.2/1.1/black!22, 4.6/1.9/black!25, 5.8/2.6/black!28, 7.4/2.2/black!25, 8.6/1.6/black!22, 10.9/1.0/black!18}
  \filldraw[fill=\c, draw=black!45] (\x,1.5) rectangle (\x+0.9,1.5+\h);
\filldraw[fill=tB!35, draw=tB!70!black] (12.4,1.5) rectangle (13.4,2.3);
\fill[black!70, rounded corners=1pt] (0.5,0) rectangle (1.1,0.9);
\fill[black!70, rounded corners=1pt] (12.9,0) rectangle (13.5,0.9);
\filldraw[fill=brown!45, draw=brown!70!black] (0,0.9) rectangle (14,1.3);
\fill[black!80] (0,1.3) rectangle (14,1.38);
\foreach \x in {0.3,0.9,...,13.8} \fill[black!35] (\x,1.42) circle (0.06);
\foreach \x/\c in {3.3/tA, 3.9/tA, 6.9/obs, 7.5/obs, 10.5/tC} \fill[\c] (\x,1.42) circle (0.07);
\fill[white, opacity=0.55] (0,1.38) rectangle (14,1.52);
\filldraw[fill=brd!40, draw=brd] (2.3,0.35) rectangle (4.3,0.85);
\node[font=\tiny] at (3.3,0.6) {ESP32};
\filldraw[fill=brd, draw=brd!60!black] (6.75,0.62) rectangle (7.25,0.88);
\filldraw[fill=adm!60, draw=adm!70!black] (1.8,0.6) circle (0.24);
\filldraw[fill=adm!60, draw=adm!70!black] (12.2,0.6) circle (0.24);
\filldraw[fill=black!15, draw=black!50] (9.3,0.35) rectangle (10.9,0.85);
\node[font=\tiny] at (10.1,0.6) {Power};}

% ---- layout helpers ----
\newenvironment{ckfig}{\begin{figure}[H]\centering}{\end{figure}}
\newcommand{\msub}[1]{\par\smallskip\noindent{\bfseries #1:}\par\nobreak\vspace{-2pt}}
\tikzset{msga/.style={draw=cli!70!black, line width=0.8pt}, cmda/.style={draw=obs, line width=1pt}, loopa/.style={draw=srv, line width=1pt}, hwa/.style={draw=black!60, line width=1.2pt, dash pattern=on 1.5pt off 1.5pt}, wanta/.style={draw=cli, line width=1.6pt}, reala/.style={draw=obs, line width=1.6pt}, siga/.style={draw=srv, line width=0.8pt}, thra/.style={draw=obs, line width=0.9pt, dash pattern=on 4pt off 2pt}, evta/.style={draw=cli, line width=1pt}}
\tikzset{iol/.style={font=\scriptsize, fill=white, inner sep=1.2pt, rounded corners=1pt, align=center}}
% legend row: start point, list of style/text
\newcommand{\ckleg}[2]{%
\coordinate (lgp) at (#1);
\foreach \ls/\lt in {#2}{%
  \draw[\ls,->] (lgp) -- ++(0.8,0) node[right, font=\scriptsize, text=black] (lgt) {\lt};
  \coordinate (lgp) at ([xshift=5mm]lgt.east);}}
% purpose table: colour (unused), runs on, technology, requirements
\newcommand{\ckinfo}[4]{\par\smallskip\noindent{\renewcommand{\arraystretch}{1.25}%
\begin{tabular}{|L{0.30}|L{0.50}|L{0.20}|}\hline
\textbf{Runs on} & \textbf{Technology} & \textbf{Requirements} \\ \hline
#2 & #3 & #4 \\ \hline
\end{tabular}}\par\smallskip}
% module data map: n in, n out, centre text, centre style, inputs, outputs, legend list
% inputs and outputs: {name}/style/{data}/arrow style
\newcommand{\iomap}[7]{%
\begin{tikzpicture}[ck]
\pgfmathsetmacro{\nmax}{max(#1,#2)}
\node[#4, line width=2pt, minimum width=30mm, minimum height=\nmax*10.5mm+2mm, font=\normalsize] (c) at (0,0) {#3};
\foreach \nm/\st/\dd/\ar [count=\ii] in {#5}{%
  \pgfmathsetmacro{\yy}{((#1+1)/2-\ii)*1.05}
  \node[\st, minimum width=25mm, minimum height=8.5mm, font=\scriptsize] (in\ii) at (-6.6,\yy) {\nm};
  \draw[\ar,->] (in\ii.east) -- node[iol]{\dd} (c.west |- in\ii);}
\foreach \nm/\st/\dd/\ar [count=\ii] in {#6}{%
  \pgfmathsetmacro{\yy}{((#2+1)/2-\ii)*1.05}
  \node[\st, minimum width=25mm, minimum height=8.5mm, font=\scriptsize] (out\ii) at (6.6,\yy) {\nm};
  \draw[\ar,->] (c.east |- out\ii) -- node[iol]{\dd} (out\ii.west);}
\node[font=\scriptsize\bfseries, text=black!60] at (-6.6,\nmax*0.525+0.4) {Inputs come from};
\node[font=\scriptsize\bfseries, text=black!60] at (6.6,\nmax*0.525+0.4) {Outputs go to};
\ckleg{-7.6,-\nmax*0.525-0.5}{#7}
\end{tikzpicture}}

\begin{document}

\thispagestyle{ckfirst}
\begin{center}
{\bfseries GEBZE TECHNICAL UNIVERSITY\par}
{Computer Engineering Department\par}
{CSE 396 Module Documentation\par}
\vspace{8pt}
{\LARGE\bfseries Çılgın Kurye\par}
\vspace{3pt}
{\large\itshape Module M4: Traffic and AI\par}
\vspace{6pt}
\textbf{Instructor:} Erkan Zergeroğlu\hspace{1.4em}\textbf{Company:} Board\&Byte\hspace{1.4em}\textbf{Group:} Group-8\hspace{1.4em}\textbf{Date:} 6 November 2026
\end{center}
\begin{center}\footnotesize\renewcommand{\arraystretch}{1.15}
\begin{tabular}{|p{0.27\linewidth}|p{0.15\linewidth}|p{0.27\linewidth}|p{0.15\linewidth}|}\hline
\textbf{Name} & \textbf{Student No.} & \textbf{Name} & \textbf{Student No.} \\ \hline
Hasan Alperen Elmas & 1801042604 & Batuhan Yanar & 220104004059 \\ \hline
Muhib Taha Boy & 230104004165 &  &  \\ \hline
\end{tabular}
\end{center}
\vspace{2pt}




% ---- table and list helpers ----
\newcolumntype{L}[1]{>{\raggedright\arraybackslash}p{\dimexpr#1\linewidth-2\tabcolsep-\arrayrulewidth\relax}}
\newcolumntype{C}[1]{>{\centering\arraybackslash}p{\dimexpr#1\linewidth-2\tabcolsep-\arrayrulewidth\relax}}
\newcommand{\ckdone}{\item[$\boxtimes$]}
\newcommand{\ckopen}{\item[$\square$]}
\newcommand{\mdot}{$\bullet$}
\newcounter{msec}
\newcounter{chap}
% ---- end of helpers ----
\newcommand{\msection}[1]{\stepcounter{msec}\section{\arabic{msec}. #1}}

\msection{Purpose}

M4 brings the city to life. It finds routes and travel times on the road network. It moves the NPC
cars that make the traffic. It also drives the AI bots that fill the empty places in a game. A bot
makes the same control message that a phone sends, so for M1 a bot and a human player look the same.

\ckinfo{srv}{Main computer inside the game server. M1 runs it once in every tick}{TypeScript, our own shortest path search (Dijkstra), a car following model for NPC cars, a state machine and path following for bots, seeded random numbers}{R1 R2 R4 R6 R11}

M4 talks only to M1, inside the same program. The figure shows what M4 gets from M1 in every tick and
what it gives back.

\begin{ckfig}
\iomap{3}{4}{\textbf{M4}\\Traffic and AI}{sv}
{{City plan file}/hwb/{intersections, roads,\\lanes, stores}/inna, {M1 Game server}/sv/{players, obstacles,\\car hits}/inna, {M1 Game server}/sv/{traffic and\\bot parameters}/inna}
{{M1 Game server}/sv/{NPC cars}/inna, {M1 Game server}/sv/{bot controls}/inna, {M1 Game server}/sv/{travel time table}/inna, {M1 Game server}/sv/{traffic metrics}/inna}
{inna/function calls inside the server program}
\end{ckfig}

\msection{Requirements covered}

M4 takes part in five requirements. The last column says what will work when M4 is finished.

\begin{longtable}[]{|L{0.06}|L{0.15}|L{0.40}|L{0.39}|}
\hline \textbf{ID} & \textbf{Requirement} & \textbf{What M4 does for it} & \textbf{When M4 is done} \\ \hline \endhead
R1 & Multiplayer game & Bots take the empty places and play like players & A game can start with only one human player \\ \hline
R2 & Delivery cycle & Travel time table from each store, which M1 uses for fair near and far points & No team gets easier delivery points than another \\ \hline
R4 & Traffic & NPC cars make real traffic on the roads & The city is full of cars that never lock up \\ \hline
R6 & Obstacles & Heavy traffic adds slow cars. An accident closes a lane. Road work closes the road for routes and cars & Cars and bots go around every obstacle \\ \hline
R11 & Engineering requirements & Talks only to M1 with plain data, tested with its own test program & M4 can be built and tested without the other modules \\ \hline
\end{longtable}

\msection{Inputs and outputs}

M1 gives the inputs at the start of each tick and takes the outputs before M3 moves the motorcycles.

\begin{longtable}[]{|L{0.11}|L{0.17}|L{0.18}|L{0.54}|}
\hline \textbf{Direction} & \textbf{Data} & \textbf{From / to} & \textbf{Content} \\ \hline \endhead
In & City plan & Shared file & Intersections, roads, lanes, stores, possible delivery points \\ \hline
In & Game data & M1 & Every tick: players, obstacles with a version number, car hit events \\ \hline
In & Parameters & M6, through M1 & Number of NPC cars, traffic effects, bot skill \\ \hline
Out & NPC cars & M1, then M3, M5, M6 & Position, direction, speed, lane, size \\ \hline
Out & Bot controls & M1, then M3 & The same record as a phone: throttle, brake, steer, deliver \\ \hline
Out & Travel times & M1 & Seconds from each team store to every possible delivery point \\ \hline
Out & Traffic metrics & M1, then M6 & Number of NPC cars, mean NPC speed, bot route changes \\ \hline
\end{longtable}

\msection{Main components}

M4 has five parts. The route planner is used by both the NPC cars and the bots.

\begin{longtable}[]{|L{0.21}|L{0.39}|L{0.40}|}
\hline \textbf{Component} & \textbf{Takes} & \textbf{Gives} \\ \hline \endhead
Road graph & City plan & 36 intersections and 60 roads in both directions, with lanes and turn costs \\ \hline
Route planner & Road graph, obstacles, traffic, parameters & Fastest routes and the travel time table, computed again when a cost changes \\ \hline
NPC traffic & Routes, players, car hits, number of cars & Moved, added and removed NPC cars \\ \hline
Bot drivers & Routes, the bot's motorcycle, game data & One control record for each bot in every tick \\ \hline
Traffic metrics & Cars and bots & Values for the admin graphs \\ \hline
\end{longtable}

\msection{How it works}

\msub{The road graph and routes}
The city has 6 roads in each direction. Where two roads cross,
there is an intersection, so the graph has 36 intersections. Between them there are 60 road
segments, and cars can drive both ways. Each segment has a cost: the expected time to drive it.
Turning left or making a U turn adds a small extra cost. Obstacles change the costs. Heavy traffic
and an accident make a segment slower, and road work closes it. M4 uses Dijkstra's shortest path
search to find the fastest routes. The graph is small, so one search is very fast. M4 computes the
routes again whenever a cost changes: when the version number of the obstacle list goes up, when the
traffic on a road gets much heavier or lighter, or when the admin changes a route parameter. From the results it makes the travel time table: the time from each of the four stores to
every possible delivery point. M1 uses this table to choose fair near and far points.

\begin{ckfig}
\begin{tikzpicture}[x=1.1cm, y=1.1cm]
\boardbase
\draw[tA, line width=2.4pt, rounded corners=3pt] (1,4) -- (1,2) -- (3,2) -- (3,1) -- (3.5,1);
\draw[tA!50, line width=1.4pt, dashed, rounded corners=3pt] (1,4) -- (3,4) -- (3,2.6);
\roadwork{3,3}{3,2}
\heavy{1,1}{2,1}
\dpoint{3.5}{1}{tA}
\foreach \x in {0,...,5}{\foreach \y in {0,...,5}{\fill[white] (\x,\y) circle (0.06);}}
\nmark{0.55}{5.15}{1}\nmark{0.6}{2.6}{2}\nmark{3.35}{3.4}{3}\nmark{3.4}{2.5}{4}\nmark{1.5}{0.65}{5}\nmark{3.95}{1.4}{6}
\node[anchor=west, font=\small, align=left] at (6.0,2.5) {\begin{tabular}{@{}c@{\ \ }l@{}}
\textbf{1} & Start: the team store\\[3pt]
\textbf{2} & Fastest route, found by Dijkstra\\[3pt]
\textbf{3} & Shorter road, but closed\\[3pt]
\textbf{4} & Road work: cost is infinite\\[3pt]
\textbf{5} & Heavy traffic: cost is higher\\[3pt]
\textbf{6} & Delivery point\\[3pt]
$\circ$ & Intersection, a node of the graph\end{tabular}};
\end{tikzpicture}
\end{ckfig}

\msub{NPC traffic}
Each NPC car follows a lane and picks a new turn at each intersection. A car
keeps a safe gap behind the car or the motorcycle in front of it. If the gap gets small, it slows
down smoothly, and it speeds up again when the road is free. A car enters an intersection only when
the intersection is free, so cars do not drive through each other. New cars never appear close to a
player. If a car is stuck for too long, M4 removes it and adds it again somewhere else, so the
traffic never locks up. The admin sets the number of cars, and heavy traffic obstacles add extra slow
cars on their road.

\begin{ckfig}
\begin{tikzpicture}[ck]
\fill[black!65] (0,0) rectangle (13,1.4);
\draw[white, dash pattern=on 6pt off 6pt] (0,0.7) -- (13,0.7);
\fill[tD!70, rounded corners=1.5pt] (1,0.15) rectangle (2.2,0.55);
\fill[adm!80, rounded corners=1.5pt] (5.2,0.15) rectangle (6.4,0.55);
\fill[tA] (9.6,0.35) circle (0.17);
\draw[cli, line width=1pt, <->] (2.3,-0.2) -- node[below, font=\scriptsize]{gap is big: speed up} (5.1,-0.2);
\draw[obs, line width=1pt, <->] (6.5,-0.2) -- node[below, font=\scriptsize]{gap is small: slow down} (9.4,-0.2);
\draw[white, ->, line width=1pt] (2.3,0.35) -- (3.3,0.35);
\draw[white, ->, line width=1pt] (6.5,0.35) -- (6.9,0.35);
\fill[black!65] (11.2,1.4) rectangle (12.6,3.0);
\fill[black!65] (11.2,-1.6) rectangle (12.6,0);
\fill[tC!80, rounded corners=1.5pt] (11.5,2.2) rectangle (11.9,2.9);
\node[font=\scriptsize, anchor=west] at (12.7,2.4) {waits until the};
\node[font=\scriptsize, anchor=west] at (12.7,2.05) {crossing is free};
\end{tikzpicture}
\end{ckfig}

\msub{Bots}
Every bot has a small state machine that follows the delivery cycle, as the figure below shows. The
bot drives to the store, gets the package, drives to the delivery point, slows down and presses
``Deliver'', and then goes back. On the way it follows the route from the planner. It steers toward a
point a little ahead on the path, so the motion looks smooth. A bot can fall or get stuck in any of
these four states. If it falls, it waits for the penalty time and then goes on in the same state. If
it is stuck, it asks for a new route and goes on in the same state. One ``bot skill'' value sets the
top speed, the reaction time and some steering noise, so the admin can make bots weaker or stronger.

\begin{ckfig}
\begin{tikzpicture}[ck]
\node[st, draw=tA, fill=tA!15] (s) at (0,0) {Go to store};
\node[st, draw=tC!80!black, fill=tC!20] (p) at (3.8,0) {Get package};
\node[st, draw=srv, fill=srvL] (d) at (7.6,0) {Drive to point};
\node[st, draw=cli, fill=cliL] (v) at (11.4,0) {Deliver};
\node[grp, fit=(s)(p)(d)(v), inner sep=9pt] (g) {};
\node[glbl, anchor=south west] at (g.north west) {ANY OF THESE STATES};
\node[st, draw=obs, fill=obs!10] (f) at (3.2,-2.6) {Fallen: wait};
\node[st, draw=adm, fill=admL] (r) at (8.2,-2.6) {Stuck: new route};
\draw[inna,->] (s) -- node[al]{arrived} (p);
\draw[inna,->] (p) -- node[al]{got it} (d);
\draw[inna,->] (d) -- node[al]{slow, near} (v);
\draw[inna,->] (v.north) -- ++(0,0.75) -| node[al, pos=0.25]{delivered, go back} (s.north);
\draw[cmda,->] ([xshift=-8mm]f.north |- g.south) -- node[al]{fall} ([xshift=-8mm]f.north);
\draw[loopa,->] ([xshift=8mm]f.north) -- node[al]{back to the same state} ([xshift=8mm]f.north |- g.south);
\draw[cmda,->] ([xshift=-8mm]r.north |- g.south) -- node[al]{stuck} ([xshift=-8mm]r.north);
\draw[loopa,->] ([xshift=8mm]r.north) -- node[al]{route ok} ([xshift=8mm]r.north |- g.south);
\ckleg{-1,-3.6}{inna/delivery cycle, cmda/problem while driving, loopa/go on in the same state}
\end{tikzpicture}
\end{ckfig}

\msection{Testing}

M4 is tested alone with a small test program that plays the role of M1. It gives the city plan, fixed
player positions and a list of test obstacles, and it draws the city from above. We check routes with
and without closed roads, equal travel times for all teams in a city that is the same on all sides,
new routes when the traffic changes, many cars in a long run without a traffic lock, bots that finish
their deliveries, and the same result with the same seed.

\end{document}
